% ===== 原文 p.12 =====
% [原文 p.12]

这些选择被编码在所谓的触手（tentacle）的一个选取之中：

\begin{enumerate}[label=（\arabic*）,leftmargin=2.2em]
\item 对每个 $j=1,\dots,4$，取 $p_j$ 处某个扇区中的一点 $a_j$，该扇区介于两条反 Stokes 射线之间\footnote{注意，触手的一个选取蕴含着对与 $p_j$ 相联系的置换给定一个标号 $P_j$（对单极点，该置换是单位矩阵）。因此，人们要么为这些选取方式固定各种约定，要么牺牲一部分几何，允许对单值性数据作循环置换。}；
\item 取基点 $p_0\in\PP^{1}/\{p_1,\dots,p_4\}$，以及对每个 $j$，一条从 $p_0$ 到 $a_j$ 的路径
\[
\gamma_j:[0,1]\to\PP^{1}/\{p_1,\dots,p_4\},
\]
使得以 $p_0$ 为基点的环路
\[
(\gamma_4^{-1}\beta_4\gamma_4)(\gamma_3^{-1}\beta_3\gamma_3)(\gamma_2^{-1}\beta_2\gamma_2)(\gamma_1^{-1}\beta_1\gamma_1)
\]
在 $\PP^{1}/\{p_1,\dots,p_4\}$ 中可缩；其中 $\beta_j$ 是 $D_j/\{p_j\}$ 中以 $a_j$ 为基点、绕 $p_j$ 正向一圈的环路。
\end{enumerate}

作为例子，考虑一个带 2 个 2 阶极点的秩 2 Higgs 丛及其在图 2.24 中的触手；其中灰色的标号与虚线路径是为定义该理论而选定的：

\begin{figure}[htbp]
\centering
\includegraphics[width=0.80\linewidth]{fig1.png}
\figtitle{2.24}{两个 2 阶极点的触手。}
\end{figure}

% [Tutor Note: 图 2.24 中灰色标号 $d_1^1,d_2^1,d_1^2,d_2^2$ 是各极点处的反 Stokes 射线（此例每个 2 阶极点恰有 2 条）；虚线中既含由 (2) 给出的从基点 $p_0$ 到 $a_j$ 的路径 $\gamma_j$，也含 $\beta_j$ 对应的绕 $p_j$ 的环路；图内文字保留原图，未重画。]

最后我们指出：若考虑点发生碰撞（collision of points）的情形，则关于野 Higgs 丛与抛物 Higgs 丛之间的关系有如下猜想：

\begin{conj}{2.25（[5]）} 对于联络中带有 $n$ 个单极点（$n$-simple poles）的奇异抛物 Higgs 丛，以及对于带单个 $n$ 阶高阶极点的野 Higgs 丛，其抛物模空间与野模空间之间存在平坦态射（flat morphism）。\end{conj}

% [Tutor Note: 猜想 2.25 首现 "$n$-simple poles"（$n$ 个单极点）；该猜想的几何图像来自本页所述的「点碰撞」（collision of points）极限：$n$ 个单极点碰撞成一个 $n$ 阶高阶极点，从而使抛物模空间与野模空间之间出现平坦态射。]
